% Section 12 — Future Work (compressed v3, full detail in Appendix E)

\section{Future Work --- The Spatiotemporal Generalisation}
\label{sec:future-work}

The framework so far runs on one axis (time). But the agents whose deployment frontier the Agentic Timeline (\S\ref{sec:agentic-timeline}) bounds inhabit \emph{both} space and time: a SWE-Bench agent navigates a codebase; a WebArena agent navigates a website; a GAIA agent navigates the open web. The Augustine Problem is the temporal face of a deeper gap whose spatial face is the \emph{Cartographic Problem}: a parallel Three Spaces ontology ($\sworld, \svisit, \sself$) and a Spatiotemporal Impossibility Theorem (Theorem~\ref{thm:sit}) generalising CIT directly. Tokens carry neither $\twall$ nor $\sworld$; token-only gradient cannot align either.

\begin{theorem}[Spatiotemporal Impossibility, SIT]
\label{thm:sit}
For any loss $\mathcal{L}(\theta) = \mathbb{E}_{x \sim \mathcal{D}}[\ell(f_\theta(x_{<i}), x_i)]$ where neither $\ell$ nor $\mathcal{D}$ contains external-metric support, $\nabla_\theta\mathcal{L}$ contains zero gradient signal aligning either $\twall$ or $\sworld$ with any internal representation.
\end{theorem}

The full development --- Three Spatial Laws (SL1/SL2/SL3, mirror of L1/L2/L3), the \emph{Agentic Frontier} hypothesis
\[
T_{\max}(A) \cdot S_{\max}(A) \;\leq\; C/\ccestST(A)
\]
(P13, pre-registered), five concrete experiments E6-E10 spanning SWE-Bench Lite, WebArena, GAIA, and MLE-Bench, and the connection to world-models literature --- appears in Appendix~\ref{app:spatiotemporal}. The implication for the autonomous-agent capability curve: \textbf{the curve saturates not because intelligence saturates, but because spatiotemporal cognition does not scale}. This is the bridge from Paper 1 to the planned Paper 3 (\emph{The Agentic Frontier}).
